% Lösungen Einheit 2 von 2 – Um den Körper herum (zusammenbau v0.1)
\einheitenkopf{Einheit 2 von 2 · Um den Körper herum}

\erg{10}{a) $f(2)$ \quad b) Der Term ist die Fläche des Querschnittskreises in der Höhe x: Kugelradius 40 (Hypotenuse), Abstand $40 - x$ des Schnitts vom Kugelmittelpunkt und Scheibenradius r bilden ein rechtwinkliges Dreieck, also $r^2 = 40^2 - (40 - x)^2$; der Faktor $\tfrac{1}{1000}$ rechnet cm³ in Liter um. \quad c) Die untere Grenze $4 - t$ zeigt: vom Boden aus wird ein Stück der Länge t erfasst; b(t) ist das Flüssigkeitsvolumen in Litern, wenn die stehende Flasche bis zur Höhe t dm gefüllt ist. \quad d) Der Rotationskörper enthält den Zylinder mit Radius 1 und Höhe 1 und den mit Radius 2 und Höhe 2; zusammen $\pi \cdot 1^2 \cdot 1 + \pi \cdot 2^2 \cdot 2 = 9\pi$, also ist sein Volumen größer – es zählen Radien, nicht Flächen unter f. \quad e) Kantenlänge der Grundfläche $= 2 \cdot 6 = 12$ cm (Durchmesser, nicht Radius; 1,2 LE = 6 cm), Höhe $= 2{,}5 \cdot 5 = 12{,}5$ cm \quad f) (1) falsch: das ist die Rotation um die x-Achse; um die y-Achse nach $x^2$ auflösen, $x^2 = 4 - 2y$, und über y integrieren: $V = \pi \cdot$ Integral von 0 bis 2 über $(4 - 2y)\,dy = 4\pi$ VE; (2) falsch: 1 VE $= 3^3 = 27$ cm³, also $V = 108\pi \approx 339{,}3$ cm³; (3) richtig, wenn die Dichte in g/cm³ angegeben ist.}
\erg{11}{a) Fehler: der Vasenradius ist der Inkreisradius des Sechsecks, nicht der Umkreisradius (nur dieser ist gleich der Seite). Richtig: $s = \tfrac{2 \cdot 1}{\sqrt{3}} \approx 1{,}15$ dm \quad b) An der Stelle des größten Funktionswerts ist der Körper am dicksten; jede Verpackung muss diesen größten Querschnittskreis umfassen. \quad c) Die Oberfläche ist ein Kreis mit dem Radius $q(2 + h)$: $A(h) = \pi \cdot (3(2 + h) - 6) = 3\pi h$; $A(4) = 12\pi \approx 37{,}7$ dm² $\approx 0{,}38$ m² – die Plane reicht}
